\subsection{Proof of Theorem \ref{t1.1}}\label{s1.2}

Let $S(TM)=S_+(TM)\oplus S_-(TM)$ be the ${\bf Z}_2$-graded Hermitian vector bundle of spinors associated to $\big(TM,g^{TM}\big)$, carrying the canonically induced Hermitian connection $\nabla^{S(TM)}=\nabla^{S_+(TM)}+\nabla^{S_-(TM)}$ (cf.~\cite{LaMi89}). And we use similar notation for $ S^n(1)$.

Following~\cite{L98}, let $E=E_+\oplus E_-$ be the ${\bf Z}_2$-graded Hermitian vector bundle
\begin{gather}\label{1.4b}
f^*\big(S\big(TS^n(1)\big)\big)=f^*\big(S_+\big(TS^n(1)\big)\big)\oplus f^*\big(S_-\big(TS^n(1)\big)\big)
\end{gather}
over $M$ carrying the pull-back Hermitian connection $\nabla^E=\nabla^{E_+}+\nabla^{E_-}$. Let $R^E=\big(\nabla^E\big)^2$ be the curvature of $\nabla^E$.

Let $D^E\colon \Gamma\big(S(TM)\widehat\otimes E\big)\rightarrow
\Gamma\big(S(TM)\widehat\otimes E\big)$ be the canonically defined (twisted by~$E$) Dirac operator (cf. \cite{LaMi89}).\footnote{Here ``$\widehat\otimes$'' is the notation for the ${\bf Z}_2$-graded tensor product (cf.~\cite{Q85}).} Then one has the canonical splitting $D=D_++D_-$ with $D^E_+\colon \Gamma(S_+(TM) \otimes E_+\oplus S_-(TM)\otimes E_-) \rightarrow \Gamma(S_-(TM) \otimes E_+\oplus S_+(TM)\otimes E_-) $, while $D^E_-\colon \Gamma(S_+(TM) \otimes E_-\oplus S_-(TM)\otimes E_+) \rightarrow \Gamma(S_-(TM) \otimes E_-\oplus S_+(TM)\otimes E_+) $. Moreover, one has formally that
\begin{gather*}%\label{1.5}
\big(D^E_+\big)^*=D^E_-.
\end{gather*}

\looseness=-1 Take any $p\in S^n(1)\setminus {f(M\setminus K)}$. Let $X\in TS^n(1)$ be a smooth vector field on $S^n(1)$ such that $|X|>0$ on $S^n(1)\setminus \{p\}$. Let $v=c(X)\colon S_+\big(TS^n(1)\big)\rightarrow S_-\big(TS^n(1)\big)$ be the Clifford action of $X$. Let $v^*\colon S_-\big(TS^n(1)\big)\rightarrow S_+\big(TS^n(1)\big)$ be the adjoint of $v$ with respect to the Hermitian metrics on $S_\pm(TS^n(1))$. Let $V\colon S\big(TS^n(1)\big)\rightarrow S\big(TS^n(1)\big)$ be the self-adjoint odd endomorphism defined by
\begin{gather}\label{1.6}
V=v+v^*.
\end{gather}
Then one has
\begin{gather}\label{1.6a}V^2=|X|^2.
\end{gather}
Thus $V$ is invertible on $ {f(M\setminus K)} = \overline {f(M\setminus K)} $. Also, $f^*V$ extends to an action on $S(TM)\widehat\otimes E$ such that for any $\alpha\in S(TM)$, $u\in E$, one has $(f^*V )\big(\alpha\widehat \otimes u\big)=(-1)^{\deg (\alpha)}\alpha\widehat\otimes (f^*V)(u)$ (cf.~\cite{Q85}).

Let $U_{\frac{1}{2}}\subset M$ be the subset defined by
$U_{\frac{1}{2}}=\big\{ x\in \operatorname{Supp}({\rm d}f)\colon |{\rm d}f(x)|<\frac{1}{2}\big\}$. Let $V_{\frac{1}{2}}\subset M$ be the open subset defined by $V_{\frac{1}{2}}=\big\{ x \colon \big|{\wedge}^2({\rm d}f(x))\big|>\frac{1}{2}\big\}$, where $\wedge^2({\rm d}f)$ is the induced action of ${\rm d}f$ on the exterior product $\wedge^2(TM)$. Clearly, $\overline{U}_{\frac{1}{2}}\cap \overline{V}_{\frac{1}{2}}=\varnothing$.

Let $\varphi\colon M\rightarrow [0,1]$ be a smooth function such that $\varphi=1$ on $(M\setminus \operatorname{Supp}({\rm d}f) )\cup U_{\frac{1}{2}}$, while $\varphi=0$ on~$V_{\frac{1}{2}}$. The existence of $\varphi$ is clear.

Similar to \cite[equation~(1.11)]{Z19}, for any $\varepsilon>0$, let $D^E_\varepsilon\colon \Gamma\big(S(TM)\widehat\otimes E\big)\rightarrow
\Gamma\big(S(TM)\widehat\otimes E\big)$ be the deformed twisted Dirac operator defined by\footnote{In view of~\cite{Q85}, one may regard $D^E_\varepsilon$ as a ``super'' Dirac operator.}
\begin{gather}\label{1.7}
D^E_\varepsilon =D^E+\varepsilon \varphi f^*V.
\end{gather}
Let $D^E_{\varepsilon,+}\colon \Gamma(S_+(TM) \otimes E_+\oplus S_-(TM)\otimes E_-) \rightarrow \Gamma(S_-(TM) \otimes E_+\oplus S_+(TM)\otimes E_-) $ and $D^E_{\varepsilon,-}\colon \Gamma(S_+(TM) \otimes E_-\oplus S_-(TM)\otimes E_+) \rightarrow \Gamma(S_-(TM) \otimes E_-\oplus S_+(TM)\otimes E_+) $ be the natural restrictions.

From (\ref{1.7}), one has
\begin{gather}\label{1.8}
\big(D^E_\varepsilon\big)^2 =\big(D^E\big)^2+\varepsilon c({\rm d}\varphi) f^*V
+\varepsilon \varphi\big[D^E,f^*V\big]+\varepsilon^2\varphi^2f^*\big(V^2\big),
\end{gather}
where $c(\cdot)$ is the notation for the Clifford action, $[\cdot,\cdot]$ is the notation for the supercommutator in the sense of~\cite{Q85}, and we identify~${\rm d}\varphi$ with the gradient of $\varphi$.

Let $e_1, \dots, e_n$ be an orthonormal basis of $\big(TM,g^{TM}\big)$. By the definition of the Dirac operator, $D^E=\sum\limits_{i=1}^n c(e_i)\nabla_{e_i}$, one has (cf.~\cite{Z19})
\begin{gather}\label{1.13}
\big[D^E,f^*V\big] =\sum_{i=1}^nc (e_i ) f^* \big(\nabla^{S(TS^n(1))}_{f_*(e_i)}V\big).
\end{gather}

Since by definition $\varphi=1$ on $M\setminus K$, while $f$ is locally constant on $M\setminus K $, from (\ref{1.8}) and~(\ref{1.13}) we see that the following identity holds on $M\setminus K$,
\begin{gather}\label{1.9}
\big(D^E_\varepsilon\big)^2 =\big(D^E\big)^2 +\varepsilon^2 f^*\big(V ^2\big).
\end{gather}

As $V|_{ {f(M\setminus K)}}$ is invertible, from (\ref{1.9}), one sees that there is a constant $a>0$ such that for any $s\in \Gamma\big(S(TM)\widehat\otimes E\big)$ supported in a~compact subset of $M\setminus K$, one has
\begin{gather}\label{1.10}
\big\|D^E_\varepsilon s\big\| \geq \varepsilon a \|s\|.
\end{gather}

 From (\ref{1.10}), one sees that one can apply the relative index theorem in~\cite{GL83} to $D^E_{\varepsilon,+}$ (compare with \cite[Section~6$\frac{4}{5}$]{Gr96}). In particular, one gets, by similar computations as in~\cite{GL83} and \cite[Proposition~III.11.24]{LaMi89},
\begin{align}
\operatorname{ind}\big(D^E_{\varepsilon,+}\big)
&= \deg (f)\big\langle {\rm ch}\big(S_+\big(TS^n(1)\big)\big) -
 {\rm ch}\big(S_-\big(TS^n(1)\big)\big),\big[S^n(1)\big]\big\rangle
\nonumber\\
& =(-1)^{\frac{n}{2}}\deg (f) \chi\big(S^n(1)\big)
 =2(-1)^{\frac{n}{2}}\deg (f) .\label{1.11}
\end{align}

 By the Lichnerowicz formula \cite{L63} for $D^E$ (cf.~\cite{LaMi89}), one has
\begin{gather}\label{1.12}
\big(D^E\big)^2=-\Delta^E+\frac{k^{TM}}{4}+\frac{1}{2}\sum_{i, j=1}^n c (e_i )c (e_j )R^E (e_i,e_j ) ,
\end{gather}
where $-\Delta^E\geq 0$ is the corresponding Bochner Laplacian.

From \cite[equation~(4.6)]{L98}, one knows that
\begin{gather}\label{1.14}
\frac{1}{2}\sum_{i, j=1}^n c (e_i )c (e_j )R^E (e_i,e_j )\geq -\frac{n(n-1)}{4}\big|{\wedge}^2({\rm d}f)\big|.
\end{gather}

From (\ref{1.1}), (\ref{1.3}), (\ref{1.8}), (\ref{1.12}) and (\ref{1.14}), one has that near any $x\in V_{\frac{1}{2}}$,
\begin{gather}\label{1.15}
\big(D_\varepsilon^E\big)^2+\Delta^E \geq0.
\end{gather}

Near any $x\in \operatorname{Supp}({\rm d}f)\setminus \overline{V}_{\frac{1}{2}}$, by
(\ref{1.3}), (\ref{1.8}), (\ref{1.12}) and (\ref{1.14}), one has
\begin{gather}\label{1.16}
\big(D_\varepsilon^E\big)^2+\Delta^E \geq \frac{n(n-1)}{8}+\varepsilon c({\rm d}\varphi)V+\varepsilon \varphi \big[D^E,f^*V\big]+\varepsilon^2\varphi^2f^*\big(V^2\big).
\end{gather}

From (\ref{1.8}), (\ref{1.13}), (\ref{1.12}) and (\ref{1.14}), one has that near any $x\in M\setminus \operatorname{Supp}({\rm d}f)$,
\begin{gather}\label{1.18}
\big(D_\varepsilon^E\big)^2+\Delta^E \geq \frac{k^{TM}}{4}+\varepsilon^2f^*\big(V^2\big).
\end{gather}

Now assume that (\ref{1.4}) does not hold. Then one has over $M$ that
 \begin{gather}\label{1.19}
k^{TM}\geq 0.
\end{gather}

From (\ref{1.13}), (\ref{1.9}), (\ref{1.15})--(\ref{1.19}) and the compactness of $\operatorname{Supp}({\rm d}f)$, we see that when $\varepsilon>0$ is small enough, there is a smooth nonnegative endormorphism $a_\varepsilon$ of $S(TM)\widehat\otimes E$ such that
 \begin{gather}\label{1.20}
a_\varepsilon>0 \qquad {\rm on}\quad (M\setminus K)\cup U_{\frac{1}{2}}
\end{gather}
and that one has on $M$ that
\begin{gather}\label{1.21}
\big(D_\varepsilon^E\big)^2 \geq -\Delta^E +a_\varepsilon.
\end{gather}

From (\ref{1.20}) and (\ref{1.21}), one finds that the relative index of $D^E_{\varepsilon,+}$ verifies that
\begin{gather*}%\label{1.22}
\operatorname{ind}\big(D^E_{\varepsilon,+}\big)=0,
\end{gather*}
which contradicts (\ref{1.2}) and (\ref{1.11}). This completes the proof of Theorem~\ref{t1.1}.
